% ===== PRÉAMBULE CANONIQUE — à coller VERBATIM en tête de CHAQUE .tex =====
\documentclass[11pt,a4paper]{article}
\usepackage[utf8]{inputenc}\usepackage[T1]{fontenc}\usepackage{lmodern}
\usepackage[french]{babel}
\usepackage{amsmath,amssymb,mathtools}\usepackage{icomma}
\usepackage[version=4]{mhchem}          % \ce{...} : équations & formules
\usepackage{siunitx}\sisetup{locale=FR} % \qty{}{}, \si{}, \ang{}
\usepackage{chemfig}                    % \chemfig{...} : structures développées
\usepackage{xcolor}\usepackage{colortbl}
\usepackage{tikz}\usepackage{pgfplots}\pgfplotsset{compat=1.18}
\usepackage[most]{tcolorbox}\usepackage{enumitem}\usepackage{array,booktabs}
\usepackage[a4paper,margin=2.2cm]{geometry}
\usetikzlibrary{arrows.meta,calc,angles,quotes,positioning,patterns,backgrounds,shapes.geometric}
\definecolor{primary}{RGB}{222,49,99}\definecolor{primarydark}{RGB}{150,22,55}
\definecolor{defcol}{RGB}{0,114,178}\definecolor{methcol}{RGB}{52,92,148}
\definecolor{excol}{RGB}{0,135,90}\definecolor{attncol}{RGB}{200,40,55}
\tcbset{boxbase/.style={enhanced,breakable,boxrule=0.9pt,arc=2.5pt,
  left=3mm,right=3mm,top=2.3mm,bottom=2.3mm,fonttitle=\bfseries,coltitle=white}}
% --- boîtes TUTORIEL (noms EXACTS, ne pas renommer) ---
\newtcolorbox{cle}[1][]{boxbase,colback=defcol!6,colframe=defcol,
  title={\textsc{À retenir}\ifx\relax#1\relax\relax\else\textmd{ : \itshape #1}\fi}}
\newtcolorbox{methode}[1][]{boxbase,colback=methcol!7,colframe=methcol,sharp corners,
  title={\textsc{Méthode}\ifx\relax#1\relax\relax\else\textmd{ --- \itshape #1}\fi}}
\newtcolorbox{exemplebox}[1][]{boxbase,colback=excol!5,colframe=excol,
  title={\textsc{Exemple}\ifx\relax#1\relax\relax\else\textmd{ : \itshape #1}\fi}}
\newtcolorbox{piege}[1][]{boxbase,colback=attncol!6,colframe=attncol,
  title={\textsc{Piège}\ifx\relax#1\relax\relax\else\textmd{ : \itshape #1}\fi}}
\newtcolorbox{remarque}[1][]{boxbase,colback=black!4,colframe=black!45,
  title={\textsc{Remarque}\ifx\relax#1\relax\relax\else\textmd{ : \itshape #1}\fi}}
% --- boîtes EXERCICE / CORRIGÉ (noms EXACTS, auto-comptées) ---
\newtcolorbox[auto counter]{exo}[1][]{boxbase,colback=primary!4,colframe=primary,
  title={\textsc{Exercice}~\thetcbcounter\ifx\relax#1\relax\relax\else\textmd{ --- \itshape #1}\fi}}
\newtcolorbox[auto counter]{corrige}[1][]{boxbase,colback=excol!4,colframe=excol,
  title={\textsc{Corrigé}~\thetcbcounter\ifx\relax#1\relax\relax\else\textmd{ --- \itshape #1}\fi}}
\newcommand{\R}{\mathbb{R}}\newcommand{\N}{\mathbb{N}}\newcommand{\Z}{\mathbb{Z}}\newcommand{\ds}{\displaystyle}
% (un \usepackage{fancyhdr}+en-tête est possible ; il est ignoré côté web)
% =========================================================================
\begin{document}
\begin{exo}[repérer la liaison dative]
Ces deux ions se forment par capture d'un proton \ce{H+}. Pour chacun, indique quelle liaison est
\textbf{dative}, quel atome est le \textbf{donneur} et lequel est l'\textbf{accepteur}.
\begin{enumerate}[label=\alph*)]
  \item \ce{NH4+} (formé à partir de \ce{NH3} $+$ \ce{H+}).
  \item \ce{H3O+} (formé à partir de \ce{H2O} $+$ \ce{H+}).
\end{enumerate}
\end{exo}

\begin{exo}[nombre total d'électrons de valence]
Calcule le nombre total d'électrons de valence de chaque édifice (attention à la charge).
\begin{enumerate}[label=\alph*)]
  \item \ce{CO2}
  \item \ce{NH4+}
  \item \ce{SO4^2-}
  \item \ce{CH4}
\end{enumerate}
\end{exo}

\begin{exo}[octet, lacune ou octet étendu ?]
Pour l'atome central souligné, dis s'il respecte l'octet, présente une \textbf{lacune}
(sous-octet) ou un \textbf{octet étendu}.
\begin{enumerate}[label=\alph*)]
  \item \ce{C} dans \ce{CH4}
  \item \ce{B} dans \ce{BF3}
  \item \ce{S} dans \ce{SF6}
  \item \ce{P} dans \ce{PCl5}
  \item \ce{Be} dans \ce{BeCl2}
\end{enumerate}
\end{exo}

\begin{exo}[qui peut dépasser l'octet ?]
Parmi ces atomes centraux, lesquels \emph{peuvent} dépasser l'octet ? Justifie par leur période
dans le tableau périodique.
\begin{enumerate}[label=\alph*)]
  \item \ce{N}
  \item \ce{S}
  \item \ce{O}
  \item \ce{Cl}
  \item \ce{C}
\end{enumerate}
\end{exo}

\begin{exo}[électronégativité et caractère ionique]
On donne les électronégativités : \ce{H}~$2{,}2$ ; \ce{Cl}~$3{,}2$ ; \ce{Na}~$0{,}9$ ;
\ce{F}~$4{,}0$. Classe les liaisons suivantes de la \textbf{moins polaire} à la \textbf{plus
polaire} (caractère ionique croissant).
\[ \ce{Cl2}\;,\quad \ce{HCl}\;,\quad \ce{NaCl}\;,\quad \ce{NaF}. \]
\end{exo}

\end{document}
